\section{Methodology}

\subsection{Hypothesis Framework}

\textbf{H-E1 (Existence):} Class-wise accuracy profiles exhibit variance beyond what overall accuracy and per-class difficulty explain.

\textbf{H-M1 (Mechanism):} Weight matrix features correlate with class-wise accuracy profiles, explaining variance beyond the stratified baseline.

\subsection{Behavioral Variance Decomposition (H-E1)}

We construct a stratified baseline predicting class-wise accuracy from overall model accuracy:
\begin{equation}
\hat{a}_{i,c} = \alpha_c \cdot A_i + \beta_c
\end{equation}

The residual variance ratio measures the fraction \emph{not} explained:
\begin{equation}
\text{Residual Ratio} = 1 - \frac{\sum_{i,c}(\hat{a}_{i,c} - a_{i,c})^2}{\sum_{i,c}(a_{i,c} - \bar{a}_c)^2}
\end{equation}

A ratio $> 0.05$ indicates meaningful behavioral variance.

\subsection{Weight Statistics Probe (H-M1)}

For each model, we extract 5 statistics (mean, std, min, max, L2 norm) from 5 layers, yielding 25 features. We use Ridge regression ($\alpha = 1.0$) to predict per-class accuracy:
\begin{equation}
\Delta R^2 = R^2_{\text{weight}} - R^2_{\text{baseline}}
\end{equation}

Success requires $\Delta R^2 > 0$.
